\documentclass{scrartcl}
\usepackage{graphicx}

% for cyrillic and Bulgarian language
\usepackage[utf8]{inputenc}
\usepackage[T2A]{fontenc}
\usepackage[bulgarian]{babel}

% math packages
\usepackage{amsmath}
\usepackage{float}
\usepackage{amsfonts}
\usepackage{amssymb}
\usepackage{amsthm}
\usepackage{cancel}

\usepackage{ragged2e}

\usepackage{listings}
\usepackage{xcolor}

\definecolor{codegreen}{rgb}{0,0.6,0}
\definecolor{codegray}{rgb}{0.5,0.5,0.5}
\definecolor{codepurple}{rgb}{0.58,0,0.82}
\definecolor{backcolour}{rgb}{0.95,0.95,0.92}

\lstdefinestyle{mystyle}{
    backgroundcolor=\color{backcolour},
    commentstyle=\color{codegreen},
    keywordstyle=\color{magenta},
    numberstyle=\tiny\color{codegray},
    stringstyle=\color{codepurple},
    basicstyle=\ttfamily\footnotesize,
    breakatwhitespace=false,
    breaklines=true,
    captionpos=b,
    keepspaces=true,
    numbers=left,
    numbersep=5pt,
    showspaces=false,
    showstringspaces=false,
    showtabs=false,
    tabsize=2,
    columns=fullflexible
}

\lstset{style=mystyle}

\usepackage[colorlinks=true, linkcolor=blue, urlcolor=cyan]{hyperref}

\newtheorem{definition}{Дефиниция}
\newtheorem{theorem}{Теорема}
\newtheorem{remark}{Забележка}

\title{XML - структуриране, валидация, обработка и представяне на документно съдържание}
\author{Ивайло Андреев + chat gpt 5.4, 5.5, gemini 3 flash}
\date{9 май 2026}

\begin{document}

\maketitle

\tableofcontents
\newpage

\begin{FlushLeft}

\section{Добре-структуриран XML}

XML (\textit{eXtensible Markup Language}) е метаезик за описание на структурирани данни и документи. Основната идея е съдържанието да се представя чрез ясно означени елементи, които са едновременно машинно обработваеми и разбираеми за човека.

\subsection{Основни концепции}

\begin{definition}
\textbf{Таг} е синтактична конструкция от вида \texttt{<name>} или \texttt{</name>}, с която се отбелязва начало или край на елемент.
\end{definition}

\begin{definition}
\textbf{Елемент} е логическа единица в XML документа, съставена от отварящ таг, съдържание и затварящ таг, например \texttt{<person>Ivan</person>}.
\end{definition}

\begin{definition}
\textbf{Атрибут} е двойка от име и стойност, записана в отварящия таг на елемента, например \texttt{id="42"}. Стойността на атрибута винаги се поставя в кавички.
\end{definition}

В началото на документа обикновено стои XML декларация, например:

\begin{lstlisting}[caption=XML declaration example]
<?xml version="1.0" encoding="UTF-8" standalone="yes"?>
\end{lstlisting}

\subsection{XML йерархии}

XML документът има дървовидна структура:

\begin{itemize}
    \item съществува точно един коренов елемент;
    \item всеки елемент може да има деца, родител и братя;
    \item елементите образуват логическа и физическа структура на документа.
\end{itemize}

Логическата структура описва отношенията между елементите, а физическата структура описва от кои единици (entities, външни ресурси) е съставен документът.

\subsection{Синтактични правила за добре-структуриран XML}

За да бъде един XML документ добре структуриран (\textit{well-formed}), трябва да са изпълнени следните условия:

\begin{itemize}
    \item всеки отварящ таг има съответен затварящ таг или е самозатварящ се;
    \item елементите са правилно вложени и не се застъпват;
    \item имената на елементите са чувствителни към регистъра;
    \item документът има точно един коренов елемент;
    \item специалните символи \texttt{<} и \texttt{\&} се екранират чрез entities, например \texttt{\&lt;} и \texttt{\&amp;}.
\end{itemize}

Важно е да се различават понятията \textit{well-formed} и \textit{valid}. Един XML документ е \textbf{well-formed}, ако спазва синтактичните правила на XML. Един XML документ е \textbf{valid}, ако освен това е well-formed и удовлетворява външна формална спецификация, например DTD или XML Schema. Това е класическо разграничение: well-formedness е синтактична коректност, а validity е съответствие със схема.

\begin{lstlisting}[caption=Well-formed XML example]
<?xml version="1.0" encoding="UTF-8"?>
<catalog>
    <book id="b1">
        <title>XML Fundamentals</title>
        <author>Alice</author>
    </book>
    <book id="b2">
        <title>XPath in Practice</title>
        <author>Bob</author>
    </book>
</catalog>
\end{lstlisting}

\subsection{Пространства от имена (Namespaces)}

\begin{definition}
\textbf{XML пространство от имена} е механизъм за избягване на конфликти между елементи и атрибути с еднакви имена, но различен смисъл.
\end{definition}

То се задава чрез атрибута \texttt{xmlns}. Често се използва префикс, който се асоциира с URI:

\begin{lstlisting}[caption=Namespace example]
<root xmlns:h="http://example.org/html"
      xmlns:f="http://example.org/furniture">
    <h:table>
        <h:tr><h:td>Cell</h:td></h:tr>
    </h:table>
    <f:table>
        <f:name>Desk</f:name>
    </f:table>
</root>
\end{lstlisting}

Префиксите като \texttt{h} и \texttt{f} са само локални имена в документа. Реалната идентичност на пространството от имена идва от URI. URI не е задължително да сочи към реален ресурс в интернет; той играе ролята на уникален идентификатор.

\section{XML валидация чрез DTD}

DTD (\textit{Document Type Definition}) описва допустимата структура на XML документ: какви елементи могат да се срещат, в какъв ред, с какви атрибути и с каква кратност.

\subsection{Цели на валидирането}

Валидацията има следните цели:

\begin{itemize}
    \item проверка дали документът следва предварително зададена структура;
    \item осигуряване на съвместим обмен на данни между системи;
    \item повторна употреба на една и съща структурна дефиниция за много документи.
\end{itemize}

Добре-структурираният документ е синтактично коректен, а валидният документ е добре-структуриран и допълнително удовлетворява дадена DTD или XML Schema.

\subsection{DTD структура и синтаксис}

DTD може да бъде:

\begin{itemize}
    \item \textbf{вътрешна} --- дефинирана в самия XML документ;
    \item \textbf{външна} --- записана в отделен файл и реферирана от XML документа.
\end{itemize}

Пример за вътрешна DTD декларация:

\begin{lstlisting}[caption=Internal DTD declaration example]
<?xml version="1.0" encoding="UTF-8"?>
<!DOCTYPE catalog [
    <!ELEMENT catalog (book+)>
    <!ELEMENT book (title, author)>
    <!ATTLIST book id ID #REQUIRED>
]>
<catalog>
    <book id="b1">
        <title>XML</title>
        <author>Alice</author>
    </book>
</catalog>
\end{lstlisting}

Пример за външна DTD декларация:

\begin{lstlisting}[caption=External DTD declaration example]
<?xml version="1.0" encoding="UTF-8"?>
<!DOCTYPE catalog SYSTEM "catalog.dtd">
<catalog>
    <book id="b1">
        <title>XML</title>
        <author>Alice</author>
    </book>
</catalog>
\end{lstlisting}

Основните DTD декларации са:

\begin{itemize}
    \item \texttt{<!ELEMENT ...>} --- описание на допустимото съдържание на елемент;
    \item \texttt{<!ATTLIST ...>} --- описание на атрибутите на елемент;
    \item \texttt{<!ENTITY ...>} --- описание на entities.
\end{itemize}

Често използвани модели на съдържание са:

\begin{itemize}
    \item \texttt{EMPTY} --- празен елемент;
    \item \texttt{ANY} --- произволно съдържание;
    \item \texttt{(\#PCDATA)} --- текстово съдържание.
\end{itemize}

Операторите за кратност са:

\begin{itemize}
    \item \texttt{?} --- нула или един път;
    \item \texttt{*} --- нула или повече пъти;
    \item \texttt{+} --- един или повече пъти.
\end{itemize}

\begin{lstlisting}[caption=DTD example]
<!ELEMENT catalog (book+)>
<!ELEMENT book (title, author, year?)>
<!ATTLIST book id ID #REQUIRED>
<!ELEMENT title (#PCDATA)>
<!ELEMENT author (#PCDATA)>
<!ELEMENT year (#PCDATA)>
\end{lstlisting}

При атрибутите често се използват типовете \texttt{CDATA}, \texttt{ID}, \texttt{IDREF} и ограниченията \texttt{\#REQUIRED}, \texttt{\#IMPLIED}, \texttt{\#FIXED}.

DTD е полезна за базово структурно валидиране, но има съществени ограничения: не поддържа богати типове данни като числа и дати и няма естествена пълна поддръжка за namespaces. Именно това е една от основните причини XML Schema да се използва като по-мощен наследник на DTD.

\section{XML валидация чрез XML Schema}

XML Schema (XSD) е по-мощен език за описание на XML структури от DTD. За разлика от DTD, самата схема е XML документ и поддържа богата типова система.

\subsection{Спецификации и типове данни}

XSD предлага:

\begin{itemize}
    \item примитивни типове като \texttt{string}, \texttt{boolean}, \texttt{decimal}, \texttt{date};
    \item производни типове като \texttt{integer}, \texttt{long}, \texttt{language}, \texttt{IDREF};
    \item потребителски типове, създадени чрез ограничаване или разширяване на съществуващи типове.
\end{itemize}

\subsection{Фасети}

\begin{definition}
\textbf{Фасети} са ограничения върху множеството допустими стойности на даден тип.
\end{definition}

Примери за фасети са:

\begin{itemize}
    \item \texttt{length}, \texttt{minLength}, \texttt{maxLength};
    \item \texttt{minInclusive}, \texttt{maxInclusive};
    \item \texttt{pattern};
    \item \texttt{enumeration}.
\end{itemize}

Кратък пример: чрез \texttt{enumeration} можем да ограничим стойността на атрибут или елемент до краен списък, например \texttt{small}, \texttt{medium}, \texttt{large}. Чрез \texttt{pattern} можем да наложим формат, например пощенски код от точно 4 цифри.

\subsection{Структури в XML Schema}

Важно е да се прави разлика между \textbf{simple content} и \textbf{complex content}. При simple content елементът съдържа само текстова стойност и няма вложени елементи. При complex content елементът може да съдържа други елементи и/или атрибути.

Основни конструкции в XSD са:

\begin{itemize}
    \item \texttt{<xs:element>} и \texttt{<xs:attribute>} --- декларации на елементи и атрибути;
    \item \texttt{<xs:simpleType>} --- тип за simple content, тоест само текстова стойност;
    \item \texttt{<xs:complexType>} --- тип за complex content, тоест вложени елементи и/или атрибути;
    \item \texttt{<xs:sequence>} --- подредена последователност;
    \item \texttt{<xs:all>} --- всички изброени елементи, без задължителен ред;
    \item \texttt{<xs:choice>} --- избор на един от няколко възможни елемента.
\end{itemize}

\begin{lstlisting}[caption=XSD example]
<xs:schema xmlns:xs="http://www.w3.org/2001/XMLSchema">
    <xs:element name="book">
        <xs:complexType>
            <xs:sequence>
                <xs:element name="title" type="xs:string"/>
                <xs:element name="author" type="xs:string"/>
                <xs:element name="year" type="xs:integer" minOccurs="0"/>
            </xs:sequence>
            <xs:attribute name="id" type="xs:ID" use="required"/>
        </xs:complexType>
    </xs:element>
</xs:schema>
\end{lstlisting}

\subsection{Сравнение между DTD и XML Schema}

\begin{center}
\begin{tabular}{|l|l|l|}
\hline
\textbf{Характеристика} & \textbf{DTD} & \textbf{XML Schema} \\
\hline
Синтаксис & Собствен синтаксис & XML синтаксис \\
\hline
Типове данни & Силно ограничени & Богата типова система \\
\hline
Namespaces & Ограничена поддръжка & Пълна поддръжка \\
\hline
Ограничения & По-базови & Фасети и сложни правила \\
\hline
\end{tabular}
\end{center}

\section{XSLT и XPath}

Технологиите XSLT и XPath се използват за избор, трансформация и представяне на XML данни.

\subsection{XSLT}

XSLT (\textit{eXtensible Stylesheet Language Transformations}) е език за трансформиране на XML документи в други формати, например HTML, текст или друг XML документ.

Обработката в XSLT е \textbf{template-based} и декларативна: описваме шаблони, които казват какви възли да бъдат разпознати и какъв резултат да се конструира, вместо да пишем императивен алгоритъм стъпка по стъпка.

Основни XSLT конструкции са:

\begin{itemize}
    \item \texttt{<xsl:template>} --- шаблон за обработка на възли;
    \item \texttt{<xsl:value-of>} --- извличане на стойност;
    \item \texttt{<xsl:for-each>} --- обхождане на множество възли;
    \item \texttt{<xsl:apply-templates>} --- прилагане на шаблони върху децата на текущия възел.
\end{itemize}

\begin{lstlisting}[caption=XSLT example]
<xsl:stylesheet version="1.0"
    xmlns:xsl="http://www.w3.org/1999/XSL/Transform">
    <xsl:template match="/catalog">
        <html>
            <body>
                <ul>
                    <xsl:for-each select="book">
                        <li><xsl:value-of select="title"/></li>
                    </xsl:for-each>
                </ul>
            </body>
        </html>
    </xsl:template>
</xsl:stylesheet>
\end{lstlisting}

XSLT позволява ясно разделяне между данни и представяне: XML пази съдържанието, а XSLT описва как то да бъде визуализирано или преобразувано.

\subsection{XPath}

\begin{definition}
\textbf{XPath} е език за адресиране и навигация в XML документ, който разглежда документа като дърво от възли.
\end{definition}

XPath използва синтаксис, подобен на файлови пътища. Основни оператори са:

\begin{itemize}
    \item \texttt{/} --- абсолютен път от корена;
    \item \texttt{//} --- търсене на възли навсякъде в документа;
    \item \texttt{@} --- избор на атрибут;
    \item \texttt{..} --- избор на родителски възел.
\end{itemize}

Особено важни са и \textbf{предикатите}, с които се филтрират възли:

\begin{itemize}
    \item \texttt{[1]} --- първият елемент от избраното множество;
    \item \texttt{[last()]} --- последният елемент;
    \item \texttt{[price > 30]} --- елементи, за които вложеният елемент \texttt{price} има стойност над 30.
\end{itemize}

Примери за XPath изрази:

\begin{itemize}
    \item \texttt{/catalog/book/title} --- всички заглавия на книги в корена \texttt{catalog};
    \item \texttt{//book[@id='b1']} --- книга с атрибут \texttt{id='b1'};
    \item \texttt{//author} --- всички елементи \texttt{author} в документа;
    \item \texttt{/catalog/book[1]} --- първата книга в каталога;
    \item \texttt{//book[last()]} --- последната книга в документа;
    \item \texttt{//book[price > 30]} --- всички книги с цена над 30.
\end{itemize}

\section{Обработка на XML: DOM и SAX}

\subsection{DOM}

DOM (\textit{Document Object Model}) е \textbf{tree-based API}: представя целия XML документ като дърво в паметта.

Основни интерфейси в DOM са:

\begin{itemize}
    \item \texttt{Node} --- базов интерфейс за възел;
    \item \texttt{Document} --- кореновият обект на документа;
    \item \texttt{Element} --- елементен възел;
    \item \texttt{Attr} --- атрибутен възел;
    \item \texttt{Text} --- текстов възел.
\end{itemize}

Предимството на DOM е, че позволява произволен достъп, навигация и модификация на документа. Недостатък е по-високата консумация на памет при големи файлове.

\subsection{SAX}

SAX (\textit{Simple API for XML}) е \textbf{event-driven streaming API}: XML документът се чете последователно, а при срещане на отварящ таг, затварящ таг, текст или грешка парсерът генерира събития.

Основни интерфейси в SAX са:

\begin{itemize}
    \item \texttt{XMLReader};
    \item \texttt{ContentHandler};
    \item \texttt{ErrorHandler}.
\end{itemize}

Приложението обработва тези събития чрез обработчици на събития като \texttt{ContentHandler} и \texttt{ErrorHandler}. SAX е особено подходящ за много големи документи, защото не съхранява цялото дърво в паметта.

\subsection{Сравнение между DOM и SAX}

\begin{center}
\begin{tabular}{|l|l|l|}
\hline
\textbf{Характеристика} & \textbf{DOM} & \textbf{SAX} \\
\hline
Модел & Дърво в паметта & Поток от събития \\
\hline
Памет & Висока консумация & Ниска консумация \\
\hline
Достъп & Произволен & Последователен \\
\hline
Модификация & Удобна & Ограничена \\
\hline
Подходящ за & Сложни промени & Много големи файлове \\
\hline
\end{tabular}
\end{center}

\section{Обобщение}

XML е универсален стандарт за структуриране и обмен на документно съдържание. Добре-структурираният XML осигурява синтактична коректност, DTD и XML Schema осигуряват валидиране спрямо формална спецификация, XSLT и XPath позволяват трансформация и навигация, а DOM и SAX дават два различни програмни модела за обработка на документи според нуждите на приложението.

\end{FlushLeft}

\end{document}
